\DocumentMetadata{
  pdfversion=2.0,
  pdfstandard=ua-2,
  lang=en-US,
  tagging=on,
  tagging-setup={math/setup=mathml-SE,tabsorder=structure}
}
\documentclass[11pt,letterpaper]{article}
\usepackage[margin=0.68in,includefoot]{geometry}
\usepackage{newcomputermodern}
\usepackage{amsmath,amscd,mathtools}
\usepackage{tikz,adjustbox}
\providecommand{\square}{\mdlgwhtsquare}
\usepackage[hidelinks]{hyperref}
\hypersetup{
  pdftitle={Analysis First-Year Exam, May 2025, Part 2},
  pdfauthor={University of Florida Department of Mathematics},
  pdfsubject={Graduate mathematics examination},
  pdfdisplaydoctitle=true,
  bookmarksopen=true
}
\setcounter{secnumdepth}{0}
\setlength{\parindent}{0pt}
\setlength{\parskip}{0.28em}
\setlength{\emergencystretch}{3em}
\setlength{\leftmargini}{2.15em}
\setlength{\leftmarginii}{2.65em}
\setlength{\leftmarginiii}{3em}
\renewcommand{\labelenumi}{\textbf{\theenumi.}}
\pagestyle{empty}
\raggedbottom
\allowdisplaybreaks
\newenvironment{examproblems}[1]{%
  \begin{enumerate}%
  \setcounter{enumi}{#1}%
  \setlength{\topsep}{0.35em}%
  \setlength{\partopsep}{0pt}%
  \setlength{\itemsep}{0.48em}%
  \setlength{\parsep}{0.18em}%
}{\end{enumerate}}
\newenvironment{examsubparts}{%
  \begin{enumerate}%
  \setlength{\topsep}{0.18em}%
  \setlength{\partopsep}{0pt}%
  \setlength{\itemsep}{0.16em}%
  \setlength{\parsep}{0.1em}%
}{\end{enumerate}}
\newenvironment{examinstructions}{%
  \par\begingroup\itshape
}{%
  \par\endgroup\vspace{0.18em}
}
\makeatletter
\renewcommand\section{\@startsection{section}{1}{0pt}%
  {-1.25ex plus -.25ex minus -.1ex}%
  {0.55ex plus .1ex}%
  {\normalfont\large\bfseries}}
\renewcommand{\@maketitle}{%
  \newpage\null\vskip -1em%
  {\footnotesize\bfseries
    \href{https://www.ufl.edu/}{University of Florida}, \href{https://math.ufl.edu/}{Department of Mathematics}\par}%
  \vskip -0.5em%
  \tagstructbegin{tag=Title}%
    {\LARGE\bfseries\href{https://gma.math.ufl.edu/exams/analysis-first-year/}{Analysis First-Year Exam}, May 2025, Part 2\par}%
  \tagstructend%
  \vskip 0.25em%
  \tagmcbegin{artifact}%
    \hrule height 1pt%
  \tagmcend%
  \vskip 1em%
}
\makeatother
\title{Analysis First-Year Exam, May 2025, Part 2}
\author{}
\date{}
\begin{document}
\maketitle
\thispagestyle{empty}
\begin{examinstructions}
Answer FOUR questions in detail. State carefully any results used without proof.
\end{examinstructions}
\begin{examproblems}{0}
\item
The sequence \((f_n)_{n=0}^{\infty}\) of real-valued functions on \([0,1]\) is uniformly bounded and converges to~\(0\) pointwise. If each \(f_n\) is (i) continuous, (ii) Riemann integrable, (iii) Lebesgue integrable, does it follow that \(\int_0^1 f_n(t)\,dt \to 0\)? In each of the three cases, a proof or counterexample should be given.
\item
Let \((f_n)_{n=0}^{\infty}\) be a sequence of continuous non-negative functions on \([0,1]\) converging pointwise to~\(0\). Prove that if the extra condition 
\[
\forall t\in[0,1]\quad f_0(t)\geq f_1(t)\geq\cdots\geq f_n(t)\geq\cdots
\]
 is satisfied then \(f_n\to 0\) uniformly on \([0,1]\). Show that this conclusion can fail if the extra condition is removed.
\item
Let the continuous real-valued function~\(f\) on \([0,1]\) satisfy 
\[
\forall n\in\mathbb N\quad \int_0^1 f(t)(1-t)^n\,dt=0.
\]
 Deduce as much as possible about~\(f\).
\item
Let \((\Omega,\mathcal A,\mu)\) be a measure space. Let \((A_n)\) be a sequence in \(\mathcal A\) and define 
\[
\liminf A_n=\bigcup_N\bigcap_{n\geq N}A_n.
\]
 Prove that 
\[
\mu(\liminf A_n)\leq\liminf\mu(A_n).
\]
\item
Let \(f(x)=\sin x/x\) for \(x>0\). Show that if~\(n\) is a positive integer then 
\[
\int_{2n\pi}^{(2n+1)\pi}\frac{\sin x}{x}\,dx\geq\frac{2}{(2n+1)\pi}.
\]
 Hence explain why~\(f\) is not Lebesgue integrable on \((a,\infty)\) for any \(a>0\).
\end{examproblems}
\end{document}
